\documentclass[ejsv2]{imsart}

%% Packages
\RequirePackage[numbers]{natbib}
%\RequirePackage[authoryear]{natbib}%% uncomment this for author-year citations
\RequirePackage[colorlinks,citecolor=blue,urlcolor=blue]{hyperref}
\RequirePackage{graphicx}

\usepackage{dsfont}

%\arxiv{2010.00000}
\startlocaldefs
%%%%%%%%%%%%%%%%%%%%%%%%%%%%%%%%%%%%%%%%%%%%%%
%%                                          %%
%% Uncomment next line to change            %%
%% the type of equation numbering           %%
%%                                          %%
%%%%%%%%%%%%%%%%%%%%%%%%%%%%%%%%%%%%%%%%%%%%%%
%\numberwithin{equation}{section}
%%%%%%%%%%%%%%%%%%%%%%%%%%%%%%%%%%%%%%%%%%%%%%
%%                                          %%
%% For Axiom, Claim, Corollary, Hypothesis, %%
%% Lemma, Theorem, Proposition              %%
%% use \theoremstyle{plain}                 %%
%%                                          %%
%%%%%%%%%%%%%%%%%%%%%%%%%%%%%%%%%%%%%%%%%%%%%%
\theoremstyle{plain}
\newtheorem{axiom}{Axiom}
\newtheorem{claim}[axiom]{Claim}
\newtheorem{theorem}{Theorem}[section]
\newtheorem{lemma}[theorem]{Lemma}
%%%%%%%%%%%%%%%%%%%%%%%%%%%%%%%%%%%%%%%%%%%%%%
%%                                          %%
%% For Assumption, Definition, Example,     %%
%% Notation, Property, Remark, Fact         %%
%% use \theoremstyle{definition}            %%
%%                                          %%
%%%%%%%%%%%%%%%%%%%%%%%%%%%%%%%%%%%%%%%%%%%%%%
\theoremstyle{definition}
\newtheorem{definition}[theorem]{Definition}
\newtheorem*{example}{Example}
\newtheorem*{fact}{Fact}
\newtheorem{assumption}{Assumption}
%%%%%%%%%%%%%%%%%%%%%%%%%%%%%%%%%%%%%%%%%%%%%%
%%                                          %%
%% For Case use \theoremstyle{remark}       %%
%%                                          %%
%%%%%%%%%%%%%%%%%%%%%%%%%%%%%%%%%%%%%%%%%%%%%%
\theoremstyle{remark}
\newtheorem{case}{Case}
%%%%%%%%%%%%%%%%%%%%%%%%%%%%%%%%%%%%%%%%%%%%%%
%% Please put your definitions here:        %%
%%%%%%%%%%%%%%%%%%%%%%%%%%%%%%%%%%%%%%%%%%%%%%
\newcommand{\E}{\mathds{E}}
\renewcommand{\P}{\mathds{P}}
\newcommand{\Z}{\mathbf{Z}}
\newcommand{\Sset}{\mathcal{S}}
\newcommand{\M}{M_\text{int}}
\newcommand{\Ps}{\mathcal{P}_\Sset}
\newcommand{\Bset}{\mathcal{B}}
\newcommand{\Pb}{\mathcal{P}_\Bset}
\newcommand{\Zn}{\Z^{(n)}}
\newcommand{\Aset}{\mathcal{A}}

\endlocaldefs

\raggedbottom
\begin{document}
\begin{frontmatter}
\title{Controlling random seed stability in regression through bagging}
%\title{A sample article title with some additional note\thanksref{t1}}
\runtitle{A sample running head title}
%\thankstext{T1}{A sample additional note to the title.}

\begin{aug}
%%%%%%%%%%%%%%%%%%%%%%%%%%%%%%%%%%%%%%%%%%%%%%%
%% Only one address is permitted per author. %%
%% Only division, organization and e-mail is %%
%% included in the address.                  %%
%% Additional information can be included in %%
%% the Acknowledgments section if necessary. %%
%% ORCID can be inserted by command:         %%
%% \orcid{0000-0000-0000-0000}               %%
%%%%%%%%%%%%%%%%%%%%%%%%%%%%%%%%%%%%%%%%%%%%%%%
\author[A]{\fnms{Nicholas}~\snm{Williams}\ead[label=e1]{nicholaswilliams@berkeley.edu}}
\and
\author[A]{\fnms{Alejandro}~\snm{Schuler}\ead[label=e2]{alejandro.schuler@berkeley.edu}}
%%%%%%%%%%%%%%%%%%%%%%%%%%%%%%%%%%%%%%%%%%%%%%
%% Addresses                                %%
%%%%%%%%%%%%%%%%%%%%%%%%%%%%%%%%%%%%%%%%%%%%%%
\address[A]{Division of Biostatistics, University of California, Berkeley\printead[presep={,\ }]{e1,e2}}
\runauthor{N. Williams et al.}
\end{aug}

\begin{abstract}

\end{abstract}

\begin{keyword}[class=MSC]
\kwd[Primary ]{00X00}
\kwd{00X00}
\kwd[; secondary ]{00X00}
\end{keyword}

\begin{keyword}
\kwd{First keyword}
\kwd{second keyword}
\end{keyword}

\end{frontmatter}
%%%%%%%%%%%%%%%%%%%%%%%%%%%%%%%%%%%%%%%%%%%%%%
%% Please use \tableofcontents for articles %%
%% with 50 pages and more                   %%
%%%%%%%%%%%%%%%%%%%%%%%%%%%%%%%%%%%%%%%%%%%%%%
%\tableofcontents

\section{Introduction}

\begin{itemize}
\item Introduce random seed stability abstractly
\item Discuss the recent increase in articles flagging random seed stability as an issue
\item Connect the issue of random seed stability to that of algorithmic stability
\item State the contributions of this work: bagging controls random seed stability
\end{itemize}

\section{Random seed stability}

\begin{itemize}
\item Discuss algorithmic stability
\item Connect algorithmic stability to random seed stability
\item Introduce the formal defintion of random seed stability
\end{itemize}

Let $\Z _i = (X_i,Y_i)$, where $X_i \in \mathcal X \subseteq \mathds{R}^d$ and, without loss of generality, $Y_i \in [0,1]$. Let $\Z = \{\Z_i\}_{i=1}^n$ denote a fixed training dataset. Let $\Sset = \left\{1, \dots, \M\right\}$ where $\M$ is the largest integer that can be represented on a machine; define $\Ps := \text{Unif}(\Sset)$. Let $\mathfrak F$ denote a generic regression algorithm and $\mathfrak{F}(\Z^{(n)}, s) = \hat f_n (\cdot; \Z^{(n)}, s) : \mathcal X \rightarrow [0,1]$ be a realized regression function, where $s \in \Sset$ is a random seed. We define the algorithm function class 
\[
\mathcal F_n (\Z) = \left\{\mathfrak F(\Z, s) : s \in \Sset \right\}.
\]
That is $\mathcal F_n (\Z)$ is the finite (albeit very large) set of functions produced by training on the same dataset and architecture, but only varying the random seed. 

\begin{definition}\label{def:stability}
An algorithm is random seed $(\epsilon,\delta)$-stable if for any training data $\Z$ and test data $x \in \mathcal X$
\[
\P_{s,s' \sim \Ps}\left(\left| \hat f_n(x;\Z,s) - \hat f_n(x;\Z,s') \geq \epsilon \right|\right) \leq \delta.
\]
\end{definition}

\begin{assumption}\label{ass:vav}
Let $\left\{\Z^{(n)}\right\}_{n=1}^\infty$ be a deterministic sequence of training datasets where $\Z^{(n)}$ has size $n$. Define
\[
\hat f_{n,\infty}(x;\Z^{(n)}) = \E_{s \sim \Ps}\left[\hat f_n(x;\Z^{(n)},s) \right].
\]
For any test data $x \in \mathcal X$, assume that as $n \rightarrow \infty$
\[
\E_{s \sim \Ps}\left[\left(\hat f_n(x;\Z^{(n)},s) - \hat f_{n,\infty}(x;\Z^{(n)}) \right)^2 \right] \rightarrow 0.
\]
\end{assumption}

\section{Subbagging}

\begin{itemize}
\item Introduce what bagging is
\item Touch on all the variants
\item Discuss how it has algorithmic stability guarantees
\item Write out the subbagging algorithm
\end{itemize}

Let
\[
B_j(m,n) = \left\{i_1, \dots, i_m \right\} \subset \left\{1, \dots, n\right\},\quad\text{where } 1 \leq m \leq n
\]
denote a $m$-size index subsample constructed without replacement. There are then $K = \binom{n}{m}$ possible subsamples; denote the collection of all the $m$-size subsamples as $\Bset = \left\{B_j(m,n) \right\}_{j=1}^K$; define $\Pb = \text{Unif}(\Bset)$.. In a slight abuse of notation, let $\Z(B_j)$ be the sub-dataset consisting of the observations of $\Z$ indexed by $B_j(m,n)$. Further, let $B_j^c(n-m,n) = \left\{1, \dots, n \right\} \setminus B_j(m,n)$ denote the index complement of $B_j(m,n)$ so that $\Z(B_j^c)$ is the sub-dataset of observations not in $\Z(B_j)$; in what follows, we will refer to these observations as \textit{out-of-bag}.

\section{Bagging controls random seed stability}

\begin{assumption}\label{ass:bvav}
Let $\left\{\Zn\right\}_{n=1}^\infty$ be a deterministic sequence of training datasets where $\Z^{(n)}$ has size $n$, and $\rho \in (0,1)$ is a fixed constant. Define $\Bset^{(n)} = \left\{B_j(m,n)\right\}_{j=1}^K$ as the set of all $K$ possible $m$-sized subsets of $\Zn$ where $m = \lfloor \rho n \rfloor$. Define
\[
\bar f_{m,\infty}(x;\Bset) = \frac{1}{K}\sum_{j=1}^K \E_{s \sim \Ps} \left[\hat f_m(x;\Z(B_j),s) \right]
\]
as the idealized subbagged estimator. For any test data $x \in \mathcal X$, assume as $n \rightarrow \infty$
\[
\E_{\zeta \sim \Ps \otimes \Pb}\left[\left(\bar f_m(x;  \Aset^{(m)},\sigma) - \bar f_{m,\infty}(x;\Bset^{(n)})\right)^2 \right] \rightarrow 0,
\]
where $\zeta = \left\{(s_1, B_1), \dots, (s_V, B_V) \right\}$ is the set of $V$ independent random seed-and-bag pairs, each drawn from the product measure $\Ps \otimes \Pb$.
\end{assumption}

\begin{theorem}\label{th:stability}
Assume $\E_{\zeta \sim \Ps \otimes \Pb}\left[\left(\bar f_m(x;  \Aset,\sigma) - \bar f_{m,\infty}(x;\Bset^{(n)})\right)^2 \right]$ decays at rate $\frac{\sigma^2}{V g(\rho n)}$. Then $\dots$ is $(\epsilon,\delta)$-random seed stable whenever $V \geq \dots $.
\end{theorem}

\subsection{Numerical experiments}

\section{Relevance to causal machine learning}

\subsection{The out-of-bag Super Learner}

\subsection{Cross-bagging}

\section{Discussion}

%%%%%%%%%%%%%%%%%%%%%%%%%%%%%%%%%%%%%%%%%%%%%%
%% Example with single Appendix:            %%
%%%%%%%%%%%%%%%%%%%%%%%%%%%%%%%%%%%%%%%%%%%%%%
\begin{appendix}
\section*{Proofs}\label{appn:proofs} %% if no title is needed, leave empty \section*{}.
\end{appendix}
%%%%%%%%%%%%%%%%%%%%%%%%%%%%%%%%%%%%%%%%%%%%%%
%% Example with multiple Appendixes:        %%
%%%%%%%%%%%%%%%%%%%%%%%%%%%%%%%%%%%%%%%%%%%%%%
%\begin{appendix}
%\end{appendix}

%%%%%%%%%%%%%%%%%%%%%%%%%%%%%%%%%%%%%%%%%%%%%%
%% Acknowledgements                         %%
%% should be provided in the                %%
%% Acknowledgements section.                %%
%%%%%%%%%%%%%%%%%%%%%%%%%%%%%%%%%%%%%%%%%%%%%%
\begin{acks}[Acknowledgments]
\end{acks}

%%%%%%%%%%%%%%%%%%%%%%%%%%%%%%%%%%%%%%%%%%%%%%
%% Funding information, if any,             %%
%% should be provided in the                %%
%% funding section.                         %%
%%%%%%%%%%%%%%%%%%%%%%%%%%%%%%%%%%%%%%%%%%%%%%
\begin{funding}
\end{funding}

%%%%%%%%%%%%%%%%%%%%%%%%%%%%%%%%%%%%%%%%%%%%%%
%% Supplementary Material, including data   %%
%% sets and code, should be provided in     %%
%% {supplement} environment with title      %%
%% and short description. It cannot be      %%
%% available exclusively as external link.  %%
%% All Supplementary Material must be       %%
%% available to the reader on Project       %%
%% Euclid with the published article.       %%
%%%%%%%%%%%%%%%%%%%%%%%%%%%%%%%%%%%%%%%%%%%%%%
\begin{supplement}
\stitle{Simulation \texttt{R} code}
\sdescription{Short description of Supplement A.}
\end{supplement}

%%%%%%%%%%%%%%%%%%%%%%%%%%%%%%%%%%%%%%%%%%%%%%%%%%%%%%%%%%%%%
%%                  The Bibliography                       %%
%%                                                         %%
%%  imsart-???.bst  will be used to                        %%
%%  create a .BBL file for submission.                     %%
%%                                                         %%
%%  Note that the displayed Bibliography will not          %%
%%  necessarily be rendered by Latex exactly as specified  %%
%%  in the online Instructions for Authors.                %%
%%                                                         %%
%%  MR numbers will be added by VTeX.                      %%
%%                                                         %%
%%  Use \cite{...} to cite references in text.             %%
%%                                                         %%
%%%%%%%%%%%%%%%%%%%%%%%%%%%%%%%%%%%%%%%%%%%%%%%%%%%%%%%%%%%%%

%% if your bibliography is in bibtex format, uncomment commands:
\bibliographystyle{imsart-number} % Style BST file (imsart-number.bst or imsart-nameyear.bst)
\bibliography{../bibliography}       % Bibliography file (usually '*.bib')

\end{document}
